\section{A principal regular-cubic CM construction and its theta indices}
\label{sec:principal-regular-twist}

The twofold construction in AC-G24 preserves a polarization of
type $(1^g,3^g)$.  Keeping the full regular cubic representation
instead permits a product principal polarization.  The integral
splitting cost and the theta-section multiplicities can both be
computed exactly.  They must be distinguished from an arithmetic
bound on the strict cohomology groups.

\begin{theorem}[AC-G27: principal regular-cubic transport]
\label{AC-G27}
Let $F\supset K=\mathbb Q(\omega)$ be a number field, and let
$A/F$ have dimension $g\geq1$ and a symmetric rigidified ample
line bundle $\mathcal L/F$ inducing a principal polarization.
Suppose $A$ has full integral CM by $\mathcal O_C$, with all
endomorphisms defined over $F$ and Rosati involution equal to
conjugation on $C$.  Let $U$ cyclically permute the three factors
of $A^3$, and twist its descent datum by the permutation
$U^{a(g)}$, where $\chi_{269}(g)=\omega^{a(g)}$.
Then the descended variety $B/F$ has the following properties.
\begin{enumerate}
\item $B$ is principally polarized by a descended line bundle
$\mathcal L_B$ whose pullback is
$\mathcal L\boxtimes\mathcal L\boxtimes\mathcal L$.
It has dimension $3g$ and stable Faltings height
$3h_{\rm Fal}^{\rm st}(A)$.
If the character is nontrivial over $F$ and $L/F$ is its cubic
splitting field, then $B\simeq\operatorname{Res}_{L/F}A_L$.
\item All elements of the CM algebra $C[U]$ are defined over
$F$.  If $C\cap K=\mathbb Q$, it is the full CM algebra
\[
                  C[U]\simeq C\times CK,\qquad
                  [C[U]:\mathbb Q]=6g.                 \tag{G27.1}
\]
Rosati conjugation is the usual conjugation on this product.
It is an etale CM algebra, not one CM field.
\item For every coefficient prime $\ell$, the Tate module is
integrally the regular Artin tensor
\[
 T_\ell B\simeq T_\ell A\otimes_{\mathbb Z_\ell}
                  \mathbb Z_\ell[C_3],                 \tag{G27.2}
\]
with diagonal Galois action and the second factor acted on by
$\chi_{269}$ through the regular permutation representation.
Writing $R_\ell=\mathcal O_C\otimes\mathbb Z_\ell$,
\[
 \operatorname{Hom}_{R_\ell}(T_\ell A,T_\ell B)
             \simeq R_\ell[C_3]                        \tag{G27.3}
\]
as integral Galois modules.  The right side has just that regular
finite-monodromy action.  This holds also at $\ell=3$.
\item If $C\cap K=\mathbb Q$, the integral CM order has the
exact sequence
\[
 0\longrightarrow\mathcal O_C[C_3]
  \longrightarrow\mathcal O_C\oplus\mathcal O_C[\omega]
  \longrightarrow\mathcal O_C/3\mathcal O_C
  \longrightarrow0.                                   \tag{G27.4}
\]
Thus its index in the displayed product order is
$3^{[C:\mathbb Q]}=3^{2g}$.
If $C$ is unramified at $3$, the second order is the maximal
order $\mathcal O_{CK}$.
\end{enumerate}
The product principal lattice is retained before this rational
splitting.  No claim is made that the larger split lattice is
again principal for the same form.
\end{theorem}
\begin{proof}
The permutation $U$ preserves the product principal polarization
integrally.  More precisely, permuting the three external tensor
factors gives canonical isomorphisms on
$\mathcal L^{\boxtimes3}$, whose third composite is the identity
and which preserve the product rigidification at zero.  This is a
genuine $C_3$-linearization.  It combines with the character
cocycle to descend the variety, line bundle, and polarization,
using the effective projective descent argument of AC-G24.
The polarization remains principal because this is checked after
base change.  Its stable height is that of $A^3$, namely three
times that of $A$.

When $L/F$ is cubic, the Weil restriction becomes $A^3$ after
base change to $L$; its descent datum permutes the three embeddings
of $L$ cyclically.  With a consistent choice of their cyclic
ordering, this is the displayed datum.  The universal property
can also be checked on an $F$-algebra $R$: the Weil restriction
has points $A(R\otimes_F L)$, and after splitting $L$ these
are the points of the three factors with precisely that action.
This proves the claimed identification.

The scalar action of $C$ commutes with $U$.  Its minimal
polynomial over $C$ on the three factors is $X^3-1$.
If $C\cap K=\mathbb Q$, its factorization is
$(X-1)(X^2+X+1)$ with the quadratic factor irreducible over $C$.
The rational Chinese remainder isomorphism gives (G27.1).
These commuting endomorphisms commute with the descent cocycle,
so they are defined over $F$.  Their total degree is twice the
dimension of $B$.  For the product polarization the adjoint of
$U$ is $U^{-1}$; on $C$ it is conjugation.  Under (G27.1)
this is exactly CM conjugation on both factors.

Applying Tate modules gives (G27.2) directly, without denominators.
A full integral $\mathcal O_C$ CM lattice is a rank-one
projective $\mathcal O_C$-module in a complex uniformization;
its $\ell$-adic completion is free of rank one over each factor
of $R_\ell$.  This is the same integral lattice argument as
in AC-G23.  Thus the $R_\ell$-linear Hom in (G27.3) is
$\operatorname{End}_{R_\ell}(T_\ell A)$ tensored with the
regular module.  The first factor is $R_\ell$ with trivial
Galois conjugation, since the CM action is defined over $F$.
The remaining action is exactly the finite regular action stated.
The weight-one representation in (G27.2) is not itself a finite
Artin representation; the weight-zero coefficient is the Hom
in (G27.3).

For (G27.4), put $R=\mathcal O_C$ and evaluate a polynomial
of degree at most two at $1$ and $\omega$.
The image consists exactly of the pairs $(a,b)$ satisfying
\[
              b\equiv a\pmod{(\omega-1)R[\omega]}.
\]
Indeed write $b=u+v\omega$.  This congruence is
$a-u-v\in3R$.  If $t=(a-u-v)/3$, the polynomial
$(u+t)+(v+t)X+tX^2$ has the prescribed evaluations.
The evaluation map is injective because a polynomial in its
kernel is a multiple $t(1+X+X^2)$ with $3t=0$, and $R$ is
$\mathbb Z$-torsion-free.  Finally
$R[\omega]/(\omega-1)\simeq R/3R$.
These facts give the exact sequence and its stated index.

If $C$ is unramified at $3$, the tensor order
$\mathcal O_C\otimes\mathcal O_K$ is maximal locally
at every rational prime.  Away from $3$ this follows because
$\mathcal O_K$ is etale over $\mathbb Z$ there, so its
base change to the normal local orders of $C$ is normal.
At $3$ the $C$-factor is etale over $\mathbb Z_3$, and
its base change to the maximal order $\mathbb Z_3[\omega]$
is again normal.  Thus the tensor order is globally normal
and equals $\mathcal O_{CK}$.  This proves the last assertion.
\end{proof}

\begin{proposition}[Theta multiplicities and the exact integral splitting index]
\label{AC-G27-theta}
For the canonical permutation linearization of AC-G27 and any
integer $d\geq1$, put
\[
 V_d=H^0(A,\mathcal L^{\otimes d}),\qquad
 r=d^g,\qquad b=(r^3-r)/3.
\]
Let $L/F$ be a finite extension splitting the descent datum
(with $L=F$ allowed for the trivial character).  Then
\[
 H^0(B,\mathcal L_B^{\otimes d})\otimes_F L
          \simeq (V_d\otimes_F L)^{\otimes_L3}.          \tag{G27.5}
\]
The multiplicities of the three eigencharacters of the geometric
automorphism $U$ are
\[
                    r+b,\quad b,\quad b.               \tag{G27.6}
\]
In particular, at level $d=1$ the unique theta section is
trivial for this canonical permutation action; a nontrivial
cubic eigenspace appears only at $d\geq2$.

The integral assertion is valid for every free module $V$ of
rank $r$ over a coefficient ring $R$: cyclic permutation of
$V^{\otimes_R3}$ has the basis-orbit decomposition
\[
 V^{\otimes_R3}\simeq R^r\oplus R[C_3]^b.             \tag{G27.7}
\]
For $R=\mathbb Z$, splitting each regular orbit into its
trivial and rational augmentation factors has index $3^b$.
For $R=\mathcal O=\mathbb Z_3[\omega]$, full Fourier
evaluation on those regular orbits has cokernel of
$\mathcal O$-length $3b=r^3-r$.
Consequently the respective logarithmic indices per module rank
are bounded by
\[
 \frac{b}{r^3}\log3\leq\frac13\log3,
 \qquad
 \frac{3b}{r^3}\log3\leq\log3.                       \tag{G27.8}
\]
The second normalization is per $\mathcal O$-rank, using
$\#(\mathcal O/\lambda)=3$; no division by
$[\mathbb Q_3(\omega):\mathbb Q_3]$ is included.
Over other integral coefficient rings the same determinants
are evaluated with their local valuations and residue norms.

For the product Hermitian metric on $V_d^{\otimes3}$, $U$ is
unitary.  Its complex character projectors are orthogonal and
have operator norm at most one.  The integral indices above
remain present despite this archimedean norm bound.
\end{proposition}
\begin{proof}
For a principal ample line bundle, abelian Riemann--Roch gives
$\dim V_d=d^g$.  An exact primary locator is J.~S. Milne,
\emph{Abelian Varieties}, version 2.00 (2008), Chapter~I,
Theorem~11.1(a)--(c), printed p.~54, in the author's
\texttt{AVc.pdf}.  To spell out the vanishing used here, the
symmetric line bundle satisfies
$[m]^*\mathcal L^{\otimes d}\simeq\mathcal L^{\otimes dm^2}$.
For large $m$ the latter is very ample and has no higher
cohomology by Theorem~11.1(c).  Pullback by the finite etale map
$[m]$, followed by its trace, is multiplication by its nonzero
degree in characteristic zero.  Therefore the higher cohomology
of $\mathcal L^{\otimes d}$ injects into zero.  The
Riemann--Roch intersection formula scales by $d^g$, and
principality gives Euler characteristic one.  This proves the
dimension formula.  The Kunneth formula for sections of an
external product gives (G27.5), compatibly with permutations.

Choose a basis $v_1,\ldots,v_r$ of $V$.  Its tensor basis is
indexed by ordered triples $(i,j,k)$.  A triple is fixed by a
three-cycle exactly when $i=j=k$, giving $r$ singleton
orbits.  Every other orbit has size three, giving $b$ regular
orbits.  This proves (G27.7).  Each regular orbit has one of
each complex character, proving (G27.6).

For a regular $\mathbb Z[C_3]$ orbit, (G27.4) with $R=\mathbb Z$
has index three.  Independent orbits multiply their indices,
giving $3^b$.  Over $\mathcal O$, the Fourier map evaluates
at $1,\omega,\omega^2$.  Its determinant is a unit times
\[
              (1-\omega)(1-\omega^2)(\omega-\omega^2).
\]
Each factor has $\lambda$-valuation one, so its cokernel
length is three.  Adding the $b$ orbit lengths gives $3b$,
while singleton orbits contribute no defect.  Division by the
rank $r^3$ proves (G27.8).

Finally the product inner product is unchanged by permuting
tensor factors.  If $U$ is unitary of order three, the operators
\[
 P_j=\frac13\sum_{a=0}^2\omega^{-ja}U^a\quad(j=0,1,2)
\]
are self-adjoint idempotents with pairwise orthogonal ranges,
as follows directly by multiplying the finite geometric sums.
An orthogonal projection has norm at most one.  The denominator
three in this formula is precisely why this analytic assertion
is compatible with the nonzero integral defect.
\end{proof}

\begin{lemma}[The unnormalized tensor-index warning]
\label{AC-G27-tensor}
If $L\subset L'$ are full $\mathbb Z$-lattices of rank $n$
and finite index $D$, then for every $d\geq1$,
\[
 [(L')^{\otimes d}:L^{\otimes d}]
            =D^{d n^{d-1}}.                             \tag{G27.9}
\]
Over a DVR, the corresponding cokernel length is
$d n^{d-1}$ times the original cokernel length.
Thus division by the new rank $n^d$ produces the normalized
cost $d/n$ times the original logarithmic index, while the
unnormalized cost need not be polynomial in $d$.
\end{lemma}
\begin{proof}
Use Smith normal form to write an inclusion matrix as
$\operatorname{diag}(a_1,\ldots,a_n)$.
The tensor matrix has entries
$a_{i_1}\cdots a_{i_d}$ on its diagonal.  In the product of
all these $n^d$ entries, each $a_i$ occurs $d n^{d-1}$
times.  Since $\prod_i|a_i|=D$, this gives (G27.9).
Taking valuations proves the DVR version.
\end{proof}

\begin{corollary}[The regular twist at the new ramified prime]
\label{AC-G27-conductor}
Relative to the product $A^3$, the regular twist has normalized
logarithmic conductor increase at most $4g\log269$.
More precisely, at $v\mid269$,
\[
 a_v(V_\ell B)\leq3a_v(V_\ell A)+4g
 \quad(\ell\ne269).                                   \tag{G27.10}
\]
If $A$ is unramified there and the cubic character is nontrivial
on inertia, the conductor is exactly $4g$, while the local Euler
factor of $B$ equals the local Euler factor of $A$.
At a good-reduction prime away from $269$, good reduction is
preserved and the Newton slopes of $B$ are those of $A$, each
repeated three times.
\end{corollary}
\begin{proof}
After extending coefficients, (G27.2) is the sum of
$V_\ell A$, $V_\ell A\otimes\chi$ and
$V_\ell A\otimes\chi^2$.
The tame-character argument of AC-G24 bounds each of the last
two conductors by $a_v(V_\ell A)+2g$, proving (G27.10).
If the original representation is unramified, the last two
summands have no inertia invariants and have conductor $2g$
each; the first retains its full invariant space.  Frobenius
on the invariant line of the regular factor is trivial, so
the local determinant is the original determinant for $A$.
Summing over $v\mid269$ as in (G24.3) proves the normalized
bound.  Away from $269$ the character is unramified.  A good
abelian scheme and its product principal polarization descend
along the unramified local splitting extension, preserving
good reduction.  Its Frobenius eigenvalues are those of $A$
multiplied by the three cubic roots prescribed by the regular
character.  Those roots are units, so their valuations, and
therefore the Newton slopes, are unchanged.
\end{proof}

\paragraph{The remaining source theorem.}
AC-G27 removes the loss of principal polarization in the
twofold scalar construction.  It supplies a genuine principal
CM object with an etale CM algebra and a canonical theta-line
linearization, together with exact integral coefficient indices.
Its algebra is a product, however, and its ramified cubic
coefficient still has the new inertia behavior at $269$.
The existing \texttt{Host} field, different-splitting ideal,
local character classes and cohomology tying data have not been
replaced by that product algebra in the supplied source.

The bounds (G27.8) concern the explicitly normalized determinant
indices of these permutation lattices.  They are not a bound
for the sizes of the strict Selmer groups, nor a comparison of
the actual differentiated theta evaluations used by the source.
To use them in its allowance one must identify the source's
normalization and prove that the relevant theta maps and integral
lattices are these descended objects or admit a controlled
comparison to them.  Equations (G27.8)--(G27.9) state the two
possible sizes exactly, rather than suppressing a dimension
factor.  The arithmetic estimate over the growing host fields
remains the separate requirement isolated in AC-G25.
